\chapter{Detection Models}
\label{chapter:model}

This chapter describes the design, training, and evaluation of the transformer-based models developed for two complementary classification tasks: binary phishing detection and multi-label emotion detection. Both tasks share the same two base architectures (\acs{RoBERTa}-large and \acs{ELECTRA}-large) and a common training philosophy, but differ in their output formulations, loss functions, and evaluation strategies. Sections~\ref{sec:spam_model_selection}--\ref{sec:spam_evaluation} cover the phishing detection models; Sections~\ref{sec:emotion_model_selection}--\ref{sec:emotion_evaluation} cover the emotion detection models.

% ─────────────────────────────────────────────────────────────────────────────
\section{Part I: Phishing Detection}
% ─────────────────────────────────────────────────────────────────────────────

\section{Model Selection}
\label{sec:spam_model_selection}

The selection of appropriate architectures for phishing email classification requires balancing representational capacity, training stability, and the specific linguistic characteristics of malicious email text. Based on a review of recent advances in \ac{NLP} and the findings surveyed in Chapter~\ref{chapter:related_work}, two transformer-based models were selected: \acs{RoBERTa}-large \cite{liu2019roberta} and \acs{ELECTRA}-large \cite{clark2020electra}.

\textbf{\acs{RoBERTa}-large} is a robustly optimised variant of \acs{BERT} trained on substantially larger data (160\,GB of text) with dynamic masking, longer sequences, and larger batch sizes. Its bidirectional self-attention mechanism allows each token to contextualise against all other tokens simultaneously, which is particularly valuable for phishing detection because manipulation tactics in malicious emails frequently span the entire message: urgency cues are planted in the subject line, credibility is built in opening paragraphs, and the malicious request arrives near the end. Capturing these long-range dependencies without vanishing gradient issues is a core strength of the RoBERTa architecture. The large variant (355\,M parameters) provides richer representational capacity than base-size alternatives, at the cost of higher memory and computational requirements, a trade-off justified here by the complexity and nuance of the detection task.

\textbf{\acs{ELECTRA}-large} departs from the masked language modelling pretraining paradigm of BERT/RoBERTa by training a discriminator network to distinguish real tokens from plausible replacements generated by a small generator model. This replaced-token-detection objective exposes the model to all input positions during pretraining, rather than only the 15\% of tokens that are masked in the BERT objective, making ELECTRA more sample-efficient and producing representations that are particularly sensitive to lexical and semantic coherence. These characteristics are directly useful for detecting subtle linguistic anomalies in phishing text: an email that is stylistically inconsistent, employs incongruous formality levels, or contains phrases that are plausible in isolation but unusual in combination will exhibit low discriminator confidence under the ELECTRA objective, providing a different signal from the contextual uncertainty captured by RoBERTa. The large discriminator (335\,M parameters) has demonstrated competitive or superior performance to RoBERTa-large on several text classification benchmarks \cite{clark2020electra}, making it a strong complementary model for ensemble construction.

The decision to use both architectures jointly was motivated by their complementary inductive biases: RoBERTa's masked language modelling produces robust long-range contextual representations, while ELECTRA's discriminative pretraining yields fine-grained sensitivity to local token plausibility. Averaging their softmax probabilities in an ensemble means that disagreements on borderline cases are naturally surfaced by the \textit{maybe spam} tier of the three-class inference scheme, rather than forcing a premature hard binary decision. Both models have demonstrated state-of-the-art performance in text classification and semantic similarity tasks \cite{liu2019roberta, sanh2019distilbert}, and their large-scale pretrained representations provide a strong initialisation for fine-tuning on the relatively modest phishing corpus assembled in Chapter~\ref{chapter:dataset}.

\section{Dataset Information \& Preparation}

The dataset used for training and evaluation consists of a combination of the original 10,000 phishing email samples generated in Chapter~\ref{chapter:dataset} and an additional 10,000 non-phishing email samples produced using the same LLM-based generation methodology, but without emotional category conditioning. The resulting corpus of 20,000 samples is balanced at the class level, with exactly 10,000 examples per class.

Before splitting, the dataset underwent a lightweight preprocessing step. The \texttt{text} column was formed by concatenating available header fields (subject, sender) with the email body, separated by newlines, to provide the model with all textual signals available in a real email client. Labels were mapped to binary integers (\texttt{non-phishing}$\to$0, \texttt{phishing}$\to$1) and rows with missing labels were discarded. Emails containing only whitespace after extraction were also removed.

Stratified splits were produced using \texttt{train\_test\_split} from scikit-learn \cite{pedregosa2011scikit} with a fixed random seed of 42 to ensure reproducibility. Twenty percent of the data was held out as the test set. A further ten percent of the remaining training data was reserved as a validation set, used exclusively for early stopping and threshold selection. This yields approximate proportions of 72\% training (14,400 samples), 8\% validation (1,600 samples), and 20\% test (4,000 samples). Stratification ensures that the 50/50 class balance is preserved across all three splits.

Class weights were computed on the training partition using \texttt{compute\_class\_weight} with the \texttt{balanced} strategy from scikit-learn. Although the overall corpus is balanced, the stratified sub-sampling can introduce minor deviations in individual partitions, so explicit class weighting serves as a safeguard and maintains consistency with the weighting strategy applied to the emotion detection models in Part~II. The resulting weights were incorporated into a custom \texttt{WeightedTrainer} that replaces the standard cross-entropy loss with a weighted variant.

\section{Model Implementation}

The implementation leveraged the Hugging Face Transformers library \cite{wolf2020transformers} for model loading, tokenization, and training, combined with scikit-learn \cite{pedregosa2011scikit} for evaluation metrics and class weight computation. For each model (\acs{RoBERTa}-large, \acs{ELECTRA}-large), the following steps were performed:

\begin{itemize}
	\item \textbf{Tokenization:} Emails were tokenised using the respective pretrained tokeniser, with padding to the maximum length and truncation at 512 tokens. Padding was handled dynamically at collation time using \texttt{DataCollatorWithPadding} to avoid inflating sequences to the global maximum during the preprocessing step, reducing memory overhead during training. For \acs{ELECTRA}-large, the \texttt{use\_fast=False} flag was set to avoid a tokeniser conversion issue present in some Transformers versions.

	\item \textbf{Class Imbalance Handling:} To address any residual class imbalance, balanced class weights were computed using scikit-learn's \texttt{compute\_class\_weight} function. These weights were incorporated into the loss function through a custom weighted trainer that replaced the standard cross-entropy loss with a weighted variant, penalizing misclassifications of the minority class more heavily. At each forward pass, the trainer pops the \texttt{labels} tensor from the inputs, computes logits from the model, and evaluates the loss with the weight tensor moved to the same device as the logits, ensuring device consistency during mixed-precision training.

	\item \textbf{Model Fine-tuning:} Each pretrained model was fine-tuned for binary sequence classification using the \ac{AdamW} optimizer \cite{loshchilov2019decoupled} with a cosine learning rate scheduler. Training was performed for up to 10 epochs with a linear warmup over the first 10\% of training steps, weight decay of 0.01, and gradient clipping with a maximum norm of 1.0. The learning rate was set to $3 \times 10^{-5}$ for \acs{RoBERTa}-large and $2 \times 10^{-5}$ for \acs{ELECTRA}-large, reflecting the greater sensitivity of the ELECTRA discriminator to higher learning rates. Label smoothing with a factor of 0.05 was applied to improve probability calibration, which is important for the threshold-based three-tier inference scheme. Early stopping with a patience of 3 epochs was used, monitoring the F1-score on the validation set, and the best model checkpoint was restored at the end of training.

	\item \textbf{Computational Optimizations:} Mixed precision training using \ac{BF16} \cite{micikevicius2018mixed} was enabled to accelerate computation and reduce memory usage on compatible \ac{GPU} hardware. Unlike the emotion detection models described in Part~II, gradient checkpointing was disabled for these binary classifiers, as sufficient \ac{GPU} memory was available and disabling it yields faster training throughput. The fused AdamW implementation (\texttt{adamw\_torch\_fused}) was used to reduce optimizer overhead, and data loading was parallelised with four worker processes and pinned memory for improved throughput.

	\item \textbf{Evaluation:} Model performance was monitored during training using accuracy, precision, recall, and F1-score on the validation set. The best model checkpoint was selected based on the highest validation F1-score. At inference time, in addition to using the standard decision threshold of 0.5, a custom threshold of 0.35 was evaluated. This lower threshold was chosen to increase recall for phishing detection, since in a security context, failing to detect a phishing email (false negative) carries greater risk than a false positive. Lowering the threshold shifts the operating point toward higher recall at a modest cost to precision, which is the preferred trade-off for a security-oriented classifier.
\end{itemize}

\section{Evaluation and Results}
\label{sec:spam_evaluation}

After training, each model was evaluated on the held-out test set. Two decision thresholds were compared: the standard threshold of 0.5, where the class with the highest softmax probability is selected, and a custom threshold of 0.35, which classifies an email as phishing whenever the predicted phishing probability exceeds 0.35. This lower threshold aims to improve recall at a modest cost to precision, reflecting the asymmetric cost of errors in phishing detection.

\subsection{ELECTRA-large Results}

Table~\ref{tab:electra_results} presents the test set performance of the \acs{ELECTRA}-large model at both decision thresholds.

\begin{table}[htbp]
\centering
\caption{Test set results for \acs{ELECTRA}-large at standard and custom thresholds.}
\label{tab:electra_results}
\begin{tabular}{lcccc}
\hline
\textbf{Threshold} & \textbf{Accuracy} & \textbf{F1} & \textbf{Precision} & \textbf{Recall} \\
\hline
0.5 (Standard) & 0.9256 & 0.9051 & 0.9230 & 0.8880 \\
0.35 (Custom)  & 0.9256 & 0.9055 & 0.9195 & 0.8920 \\
\hline
\end{tabular}
\end{table}

The \acs{ELECTRA}-large model achieved an F1-score of 0.9051 at the standard threshold, with high precision (0.9230) but comparatively lower recall (0.8880). Lowering the threshold to 0.35 yielded a marginal improvement across all metrics, increasing the F1-score to 0.9055 and recall to 0.8920, while precision decreased only slightly to 0.9195. This confirms that the custom threshold provides a better trade-off for phishing detection, capturing more true positives with minimal impact on precision.

\subsection{RoBERTa-large Results}

Table~\ref{tab:roberta_results} presents the test set performance of the \acs{RoBERTa}-large model at both decision thresholds.

\begin{table}[htbp]
\centering
\caption{Test set results for \acs{RoBERTa}-large at standard and custom thresholds.}
\label{tab:roberta_results}
\begin{tabular}{lcccc}
\hline
\textbf{Threshold} & \textbf{Accuracy} & \textbf{F1} & \textbf{Precision} & \textbf{Recall} \\
\hline
0.5 (Standard) & 0.9352 & 0.9169 & 0.9410 & 0.8940 \\
0.35 (Custom)  & 0.9336 & 0.9150 & 0.9371 & 0.8940 \\
\hline
\end{tabular}
\end{table}

The \acs{RoBERTa}-large model achieved an F1-score of 0.9169 at the standard threshold, with very high precision (0.9410) and recall (0.8940). Lowering the threshold to 0.35 resulted in a slight decrease in accuracy and F1-score, with precision dropping to 0.9371 while recall remained unchanged at 0.8940. This suggests that for \acs{RoBERTa}-large, the standard threshold already provides a strong balance between precision and recall, and the custom threshold does not offer a significant improvement in this case.

\subsection{Baseline Comparison}
\label{sec:baseline_comparison}

To contextualise the performance of the fine-tuned transformer models, two classical baselines were evaluated under the same experimental conditions: identical data splits (random seed 42), identical stratification strategy, and identical threshold settings.

\textbf{TF-IDF + Logistic Regression} represents the standard strong baseline for text classification. Input emails were tokenised into unigrams and bigrams, weighted by sublinear TF-IDF, and classified using an $\ell_2$-regularised logistic regression with balanced class weights, mirroring the class-weight strategy applied during transformer fine-tuning.

\textbf{TF-IDF + Multinomial Naive Bayes} provides a historically relevant reference point, reflecting the generation of spam filters from which phishing detection evolved \cite{metsis2006spam}. Both classifiers used a vocabulary of up to 100,000 features with a minimum document frequency of two.

Table~\ref{tab:baseline_comparison} reports all four models at both evaluated thresholds.

\begin{table}[htbp]
\centering
\caption[Phishing detection results on the held-out test set]{Phishing detection results on the held-out test set for classical baselines and fine-tuned transformer models.}
\label{tab:baseline_comparison}
\resizebox{\textwidth}{!}{%
\begin{tabular}{llcccc}
\hline
\textbf{Model} & \textbf{Threshold} & \textbf{Accuracy} & \textbf{F1} & \textbf{Precision} & \textbf{Recall} \\
\hline
\multicolumn{6}{l}{\textit{Classical baselines}} \\
TF-IDF + Multinomial NB     & 0.50 (Standard) & 0.8000 & 0.7459 & 0.7583 & 0.7340 \\
TF-IDF + Logistic Regression & 0.50 (Standard) & 0.8856 & 0.8560 & 0.8621 & 0.8500 \\
\hline
\multicolumn{6}{l}{\textit{Fine-tuned transformer models}} \\
\acs{ELECTRA}-large          & 0.35 (Custom)   & 0.9256 & 0.9055 & 0.9195 & 0.8920 \\
\acs{RoBERTa}-large          & 0.35 (Custom)   & 0.9336 & 0.9150 & 0.9371 & 0.8940 \\
\hline
\end{tabular}%
}
\end{table}

The TF-IDF + Logistic Regression baseline achieved an F1-score of 0.8560, demonstrating that lexical features alone carry substantial signal for phishing detection. However, this represents a gap of approximately 6 percentage points in F1 relative to \acs{RoBERTa}-large (0.9150), confirming that contextual, deep representations learned during pre-training provide meaningful improvements beyond surface-level term statistics. Multinomial Naive Bayes performs considerably worse (F1 = 0.7459), consistent with its inability to model feature correlations and its poor calibration on long, variable-length email texts. The ROC-AUC of the logistic regression baseline (0.948) is noteworthy --- it approaches transformer-level discriminative ability at the ranking level --- yet this does not translate into calibrated decisions at any single threshold, which explains the F1 gap. Taken together, these results establish that the transformer ensemble provides a clear and quantifiable improvement over classical approaches, while the logistic regression baseline sets a competitive non-trivial floor that any proposed method should be expected to surpass.

\subsection{Comparative Analysis}
\label{sec:evaluation}

Table~\ref{tab:phishing_model_comparison} summarizes both models using the threshold configuration selected for deployment.

\begin{table}[htbp]
\centering
\caption{Cross-model comparison on the test set (selected threshold per model).}
\label{tab:phishing_model_comparison}
\resizebox{\textwidth}{!}{%
\begin{tabular}{llcccc}
\hline
\textbf{Model} & \textbf{Selected Threshold} & \textbf{Accuracy} & \textbf{F1} & \textbf{Precision} & \textbf{Recall} \\
\hline
\multicolumn{6}{l}{\textit{Classical baselines}} \\
TF-IDF + Multinomial NB      & 0.50 (Standard) & 0.8000 & 0.7459 & 0.7583 & 0.7340 \\
TF-IDF + Logistic Regression & 0.50 (Standard) & 0.8856 & 0.8560 & 0.8621 & 0.8500 \\
\hline
\multicolumn{6}{l}{\textit{Fine-tuned transformer models}} \\
\acs{ELECTRA}-large           & 0.35 (Custom)   & 0.9256 & 0.9055 & 0.9230 & 0.8880 \\
\acs{RoBERTa}-large           & 0.35 (Custom)   & 0.9336 & 0.9150 & 0.9371 & 0.8940 \\
\hline
\end{tabular}%
}
\end{table}

Based on the similar performance of both models, the final selection was a combination of \acs{ELECTRA}-large and \acs{RoBERTa}-large, leveraging the strengths of both architectures. The ensemble approach averages the softmax probabilities of both models and applies a tiered classification scheme: emails with an averaged probability below the ensemble threshold are classified as \textit{ham}; emails between the ensemble threshold and 0.50 are flagged as \textit{maybe spam}, surfacing borderline cases for user review without a hard binary decision; and emails at or above 0.50 are classified as \textit{spam}. This three-tier scheme reflects the asymmetric risk of false negatives in a security context while avoiding the false-positive burden of a purely low-threshold binary classifier.


% ─────────────────────────────────────────────────────────────────────────────
\section{Part II: Emotion Detection}
% ─────────────────────────────────────────────────────────────────────────────

\section{Model Selection}
\label{sec:emotion_model_selection}

To model the emotional content present in phishing emails, the same two transformer-based architectures used for phishing detection were employed: \acs{RoBERTa}-large \cite{liu2019roberta} and \acs{ELECTRA}-large \cite{clark2020electra}. These models are widely adopted in \acs{NLP} classification tasks and provide strong contextual representations that are suitable for capturing subtle emotional cues in text.

Unlike the binary phishing detection task in Part~I, emotion detection is formulated here as a \textbf{multi-label classification} problem: a single phishing email may simultaneously express multiple emotions (e.g., fear and curiosity). This formulation better reflects realistic phishing language, where overlapping emotional manipulation strategies are common.

\section{Emotion Taxonomy Consolidation}
\label{sec:emotion_taxonomy}

The original annotation schema comprised fourteen emotional categories. Prior to model training, an inter-label co-occurrence analysis was conducted on the annotated corpus to identify clusters of emotions that consistently co-occurred and represented semantically equivalent manipulation vectors. This analysis motivated a principled consolidation from fourteen classes to eight, with the goal of increasing the per-class training signal while preserving meaningful distinctions between manipulation strategies.

Table~\ref{tab:emotion_merging} summarises the mapping from original categories to the consolidated taxonomy.

\begin{table}[htbp]
\centering
\caption[Emotion label consolidation]{Emotion label consolidation: mapping from the original 14-class annotation schema to the 8-class training taxonomy.}
\label{tab:emotion_merging}
\begin{tabular}{lll}
\hline
\textbf{Consolidated Label} & \textbf{Original Categories} & \textbf{Rationale} \\
\hline
\texttt{positive\_arousal} & excitement, surprise, joy, desire & Co-occurrence rate 43–60\%; \\
                           &                                    & shared positive-affect manipulation vector \\
\texttt{warmth}            & admiration, gratitude, caring      & Trust-building emotional cluster \\
\texttt{threat}            & anger, fear                        & Intimidation/coercion cluster; \\
                           &                                    & 79\% of anger samples contain fear \\
\texttt{curiosity}         & curiosity                          & Distinct manipulation tactic \\
\texttt{confusion}         & confusion                          & Distinct manipulation tactic \\
\texttt{sadness}           & sadness                            & Distinct manipulation tactic \\
\texttt{relief}            & relief                             & Distinct manipulation tactic \\
\hline
\multicolumn{3}{l}{\textit{Dropped:} \texttt{neutral} (143 samples) and \texttt{unsure} (119 samples) — too rare} \\
\multicolumn{3}{l}{to learn from; \texttt{unsure} is an annotation artefact rather than a genuine emotion.} \\
\hline
\end{tabular}
\end{table}

The three merged clusters were motivated by concrete empirical and semantic criteria. The \texttt{positive\_arousal} cluster groups emotions that all serve the same attacker goal: inducing a positive affective state in the victim to lower critical thinking and prompt impulsive action. Co-occurrence rates between any two members of this cluster ranged from 43\% to 60\%, indicating that annotators overwhelmingly labelled these emotions together rather than in isolation. The \texttt{warmth} cluster consolidates the trust-building manipulation tactics that phishing attackers use to establish rapport before delivering a malicious request; admiration, gratitude, and caring rarely appeared without one another in the corpus. The \texttt{threat} cluster is motivated by the strongest empirical signal: 79\% of samples labelled with anger also carried a fear label, and both emotions serve the same intimidation and coercion function within phishing messages.

Four categories were retained as independent classes because each represents a genuinely distinct manipulation tactic with a unique psychological profile and no significant co-occurrence with other classes: \texttt{curiosity} (exploiting information-seeking behavior), \texttt{confusion} (disorienting the victim to reduce scrutiny), \texttt{sadness} (guilt-induction and empathy exploitation), and \texttt{relief} (offering resolution of a fabricated threat to prompt compliance).

The \texttt{neutral} and \texttt{unsure} labels were dropped entirely. Their combined sample count (262 instances) was insufficient to provide reliable training signal, and \texttt{unsure} is an annotation artefact that reflects annotator uncertainty rather than a phishing manipulation strategy.

This consolidation yields a final training taxonomy of \textbf{eight classes}. Each class benefits from two to four times the training signal of its constituent original categories, which is expected to improve model calibration and reduce the severity of the class imbalance problem relative to the original schema.

The label merging was applied programmatically prior to training via a dedicated preprocessing script (\texttt{merge\_labels.py}), which reads the raw annotated corpus and writes the consolidated version to \texttt{combined\_emails\_dataset\_merged.csv}. All training scripts operate on this merged file.

\section{Dataset Preparation}

The starting point was the 20,000 combined email dataset used throughout this thesis. For emotion modelling, only emails labelled as \texttt{phishing} were retained, since the objective is to characterise emotional patterns specifically within malicious messages.

Rows without emotion annotation were removed, and the remaining emotion labels were parsed from a delimiter-based format (e.g., \texttt{curiosity\#threat\#warmth}) into sets of class tags reflecting the consolidated taxonomy described in Section~\ref{sec:emotion_taxonomy}. A global emotion vocabulary of eight classes was built from the training corpus, and each email was encoded as a multi-hot vector, where each dimension indicates the presence (1) or absence (0) of a given consolidated emotion.

To preserve label distribution during splitting, a stratification proxy was used: for each sample, the index of its dominant label (\texttt{argmax} over the multi-hot vector) served as the stratification target. Data was split into test (20\%) and train (80\%), followed by a validation split (10\% of train), resulting in approximately 72\%/8\%/20\% train/validation/test partitions, with a fixed random seed for reproducibility.

\section{Model Implementation}

Model development used the Hugging Face Transformers ecosystem \cite{wolf2020transformers} for tokenization, model fine-tuning, and checkpoint management, together with scikit-learn \cite{pedregosa2011scikit} for metric computation.

For each model (\acs{RoBERTa}-large and \acs{ELECTRA}-large), the following pipeline was applied:

\begin{itemize}
	\item \textbf{Tokenization:} Emails were tokenised with each model's pretrained tokeniser, padded to fixed length, and truncated at 512 tokens.

	\item \textbf{Multi-label Objective:} The classifier head was configured with \texttt{problem\_type=multi\_label\_classification}. During training, logits were optimized using \texttt{BCEWithLogitsLoss}, and probabilities were obtained at evaluation time using the sigmoid function.

	\item \textbf{Class Imbalance Handling:} Because some consolidated classes are less frequent than others, per-class positive weights were computed from the training split as the negative-to-positive ratio. To avoid excessive penalisation of rare classes, a square-root dampening strategy was used and values were clipped to a maximum of 5.0. These \texttt{pos\_weight} values were injected into \texttt{BCEWithLogitsLoss} through a custom weighted trainer.

	\item \textbf{Model Fine-tuning:} All models were fine-tuned for up to 10 epochs using \acs{AdamW} \cite{loshchilov2019decoupled}, cosine learning-rate scheduling, learning rate $2 \times 10^{-5}$, warmup ratio 0.15, and weight decay 0.01. No label smoothing was applied (factor 0.0), as the multi-label sigmoid objective differs from the cross-entropy setting used for the binary classifiers in Part~I. Early stopping (patience = 3) monitored validation macro-F1, with best-checkpoint restoration enabled.

	\item \textbf{Computational Optimisations:} \ac{BF16} mixed precision \cite{micikevicius2018mixed} and gradient clipping (max norm 1.0) were used. Unlike the binary spam classifiers, gradient checkpointing was enabled here to accommodate the larger per-sample memory footprint of multi-label training with a 32-sample batch size per device. This trades additional compute for reduced peak \ac{GPU} memory consumption.

	\item \textbf{Threshold Tuning:} Since multi-label outputs depend strongly on decision thresholds, threshold optimisation was performed on the validation set after training. Two strategies were evaluated: (i) one global threshold shared by all classes, and (ii) per-class thresholds. Candidate thresholds in the interval [0.10, 0.80] with 0.05 steps were tested to maximise macro-F1.

	\item \textbf{Evaluation:} Performance was reported using exact-match accuracy, macro precision, macro recall, and macro F1, along with per-class classification reports on the held-out test set.
\end{itemize}

\section{Evaluation and Results}
\label{sec:emotion_evaluation}

After training, each model was evaluated on the test set using both thresholding strategies selected from validation: global threshold and per-class thresholds. This comparison assesses whether class-specific calibration improves the balance between precision and recall across the eight consolidated emotion classes.

\subsection{RoBERTa-large Results}

Table~\ref{tab:roberta_emotion_results} summarises the test performance of the \acs{RoBERTa}-large model under both thresholding strategies.

\begin{table}[htbp]
\centering
\small
\caption{Test set results for \acs{RoBERTa}-large using global and per-class thresholds.}
\label{tab:roberta_emotion_results}
\resizebox{\textwidth}{!}{%
\begin{tabular}{lcccc}
\hline
\textbf{Threshold Strategy} & \textbf{Exact Match} & \textbf{Macro F1} & \textbf{Macro Precision} & \textbf{Macro Recall} \\
\hline
Global threshold (best on val)     & 0.2918 & 0.6869 & 0.6823 & 0.7091 \\
Per-class thresholds (best on val) & 0.2816 & 0.7041 & 0.6793 & 0.7477 \\
\hline
\end{tabular}
}
\end{table}

\subsection{ELECTRA-large Results}

Table~\ref{tab:electra_emotion_results} reports the corresponding results for \acs{ELECTRA}-large.

\begin{table}[htbp]
\centering
\small
\caption{Test set results for \acs{ELECTRA}-large using global and per-class thresholds.}
\label{tab:electra_emotion_results}
\resizebox{\textwidth}{!}{%
\begin{tabular}{lcccc}
\hline
\textbf{Threshold Strategy} & \textbf{Exact Match} & \textbf{Macro F1} & \textbf{Macro Precision} & \textbf{Macro Recall} \\
\hline
Global threshold     & 0.3326 & 0.6967 & 0.7168 & 0.6911 \\
Per-class thresholds & 0.2775 & 0.7036 & 0.6890 & 0.7407 \\
\hline
\end{tabular}
}
\end{table}

\subsection{Comparative Analysis}

Table~\ref{tab:emotion_model_comparison} provides a cross-model summary using the best threshold strategy for each architecture. The label consolidation described in Section~\ref{sec:emotion_taxonomy} ensures that all models are trained and evaluated against the same eight semantically coherent classes, making cross-model comparison meaningful and directly comparable.

\begin{table}[htbp]
\centering
\small
\caption{Cross-model comparison on the test set (best threshold strategy per model).}
\label{tab:emotion_model_comparison}
\resizebox{\textwidth}{!}{%
\begin{tabular}{lccccc}
\hline
\textbf{Model} & \textbf{Best Strategy} & \textbf{Exact Match} & \textbf{Macro F1} & \textbf{Macro Precision} & \textbf{Macro Recall} \\
\hline
\acs{RoBERTa}-large   & Per-class thresholds & 0.2816 & 0.7041 & 0.6793 & 0.7477 \\
\acs{ELECTRA}-large   & Global threshold     & 0.3326 & 0.6967 & 0.7168 & 0.6911 \\
\hline
\end{tabular}
}
\end{table}